\documentclass[aspectratio=169,10pt,spanish]{beamer}

\input{../../../_preambulo_beamer.tex}
\input{../../../_comandos_beamer.tex}

\selectlanguage{spanish}

\title{MA1001B - Modelación Estadística para la Toma de Decisiones}
\subtitle{Lección 12: Sesgo de muestreo, calidad de datos y error no muestral}
\author{}
\institute{Tecnol\'ogico de Monterrey}
\date{}

\begin{document}

\begin{frame}[plain]
  \vspace{1.2cm}
  {\Large\bfseries MA1001B - Modelación Estadística para la Toma de Decisiones\par}
  \vspace{0.7cm}
  {\large Lección 12: Sesgo de muestreo, calidad de datos y error no muestral\par}
  \vspace{0.7cm}
  {\color{TecRojo}\rule{\textwidth}{1.2pt}}
  \vspace{0.8cm}

  Facultad MA1001B\\[0.3cm]
  Periodo\\[0.3cm]
  Tecnol\'ogico de Monterrey
\end{frame}

\begin{frame}{Cuando el problema es el marco}
  \begin{alertblock}{Pregunta guía}
    ¿Puede una muestra más grande, o una tabla limpia, reparar una
    estimación de tasa de quejas si falta el canal fácil o toda una región?
  \end{alertblock}

  \vspace{0.6em}
  Al final de esta lección, usted puede:
  \begin{itemize}
    \item estimar una tasa poblacional a partir de una muestra aleatoria simple;
    \item contrastar ese error de muestreo con el sesgo de conveniencia solo en línea;
    \item identificar un hueco de cobertura cuando una región está ausente del marco;
    \item explicar por qué limpiar filas incompletas no restaura unidades faltantes.
  \end{itemize}

  \vfill
  \scriptsize Todos los tickets, regiones y banderas de queja usados hoy son sintéticos. No se usan datos reales de empresa.
\end{frame}

\begin{frame}{Un registro sintético de 8{,}000 tickets}
  \small
  \begin{center}
    \begin{tabular}{lrrr}
      \toprule
      \textbf{Región} & \textbf{Tickets} & \textbf{En el marco MAS} & \textbf{En el marco sin Oeste} \\
      \midrule
      Norte & 2{,}400 & sí & sí \\
      Sur & 2{,}400 & sí & sí \\
      Este & 2{,}000 & sí & sí \\
      Oeste & 1{,}200 & sí & no \\
      \midrule
      Total & 8{,}000 & 8{,}000 & 6{,}800 \\
      \bottomrule
    \end{tabular}
  \end{center}

  \vspace{0.4em}
  \begin{exampleblock}{Parámetro}
    La tasa poblacional de quejas es la media de la bandera completa de
    queja, incluido el Oeste. El error de muestreo no es lo mismo que un
    hueco de cobertura.
  \end{exampleblock}
\end{frame}

\begin{frame}[fragile]{Paso 1: Muestra aleatoria simple}
  \begin{columns}[T]
    \begin{column}{0.42\textwidth}
      \small
      \begin{lstlisting}
srs = register.sample(
    n=400,
    replace=False,
    random_state=42,
)
      \end{lstlisting}

      \begin{block}{Salida verificada}
        $N=8{,}000$\\
        Tasa poblacional $=0.119875$\\
        MAS $n=400$ tasa $=0.105000$\\
        Error absoluto $=0.014875$
      \end{block}
    \end{column}
    \begin{column}{0.58\textwidth}
      \centering
      \includegraphics[width=\textwidth]{../figuras/sesgo_muestreo_01_mas.png}
      \vspace{0.2em}
      \scriptsize Parámetro poblacional versus MAS del Paso 1
    \end{column}
  \end{columns}
\end{frame}

\begin{frame}{Los tickets en línea son un canal fácil y sesgado}
  \small
  Los tickets en línea se quejan con menor frecuencia que los de tienda:
  \[
    \hat{p}_{\text{en línea}}=0.057323,\qquad
    \hat{p}_{\text{tienda}}=0.176555.
  \]
  Una muestra de conveniencia de los primeros $1{,}200$ tickets en línea es
  más grande que la MAS, pero copia la tasa del canal En línea en lugar de
  la tasa poblacional.
\end{frame}

\begin{frame}[fragile]{Paso 2: Muestra de conveniencia en línea}
  \begin{columns}[T]
    \begin{column}{0.40\textwidth}
      \small
      \begin{lstlisting}
web = register.loc[
    register["channel"] == "En linea"
]
conv = web.head(1200)
      \end{lstlisting}

      \begin{block}{Salida verificada}
        Conveniencia $n=1{,}200$\\
        Tasa $=0.056667$\\
        Error $=0.063208$\\
        Más grande, más error
      \end{block}
    \end{column}
    \begin{column}{0.60\textwidth}
      \centering
      \includegraphics[width=\textwidth]{../figuras/sesgo_muestreo_02_conveniencia_web.png}
      \vspace{0.2em}
      \scriptsize Tamaño de muestra versus error del Paso 2
    \end{column}
  \end{columns}
\end{frame}

\begin{frame}{Paso 3: La limpieza no restaura unidades faltantes}
  \begin{columns}[T]
    \begin{column}{0.42\textwidth}
      \small
      Tasa de quejas del Oeste $=0.276667$. Omitir el Oeste produce error
      $0.027669$.

      \vspace{0.4em}
      Tras la eliminación listwise de banderas faltantes: $6{,}437$ filas,
      error $0.026509$, tickets del Oeste $=0$.

      \vspace{0.4em}
      \textbf{Límite:} limpiar un marco sesgado no restaura unidades
      faltantes. La Lección 13 modela una variable aleatoria discreta en un
      espacio muestral completo.
    \end{column}
    \begin{column}{0.58\textwidth}
      \centering
      \includegraphics[width=\textwidth]{../figuras/sesgo_muestreo_03_hueco_cobertura.png}
      \vspace{0.2em}
      \scriptsize Hueco de cobertura y limpieza del Paso 3
    \end{column}
  \end{columns}
\end{frame}

\begin{frame}{Verificación de conceptos}
  \begin{enumerate}
    \item ¿Cuál es el parámetro poblacional en este registro?
    \item ¿Por qué $n=1{,}200$ puede producir más error que $n=400$?
    \item ¿Excluir el Oeste es un error de muestreo o un error de cobertura?
    \item Tras la limpieza, ¿por qué el conteo del Oeste sigue siendo $0$?
  \end{enumerate}

  \vspace{0.6em}
  \begin{block}{Conexión con la Lección 13}
    La sesión puede continuar directamente con variables aleatorias
    discretas, FMP, FDA, esperanza y varianza.
  \end{block}

  \vfill
  \scriptsize Fuente: OpenStax, \textit{Introductory Business Statistics 2e},
  Secciones 1.2 y 1.4. CC BY-NC-SA.
\end{frame}

\appendix



\end{document}
